\documentclass[11pt]{article}

\usepackage[margin=1in]{geometry}
\usepackage{booktabs}
\usepackage{graphicx}
\usepackage{amsmath}
\usepackage{amssymb}
\usepackage{url}
\usepackage[hidelinks]{hyperref}
\usepackage{microtype}

% Dirac notation without pulling in the physics package, which is not always
% present in a minimal TeX installation.
\newcommand{\ket}[1]{\left|#1\right\rangle}

\title{AegisQ-HPC: Communication-Aware Distributed Quantum Simulation\\
       with Post-Quantum Secure Job Provenance}
\author{Pinar Aksoy}
\date{Technical report, \today}

\begin{document}
\maketitle

\begin{abstract}
Simulating an $n$-qubit quantum circuit costs $2^n$ amplitudes, so beyond
roughly thirty qubits the state vector must be partitioned across the memory
of several machines. Partitioning turns some quantum gates into network
operations, and which gates depends entirely on an arbitrary choice: which
logical qubits are assigned to the bits of the basis index that select a rank.
We present AegisQ-HPC, a distributed state-vector simulator that makes this
choice deliberately, and pairs it with a gate-fusion pass that reduces how
often the remaining communicating gates fire. Every byte the runtime hands to MPI is instrumented, an
analytical cost model predicts that traffic from the circuit structure, and a
search over qubit placements minimises it. On a single 8-core host with
shared-memory MPI, communication-aware placement removes $87.2\%$ of measured
MPI traffic for a 20-qubit quantum Fourier transform on 8 ranks (wall time
$-24.3\%$) and $72.3\%$ for a Grover circuit, while leaving a GHZ chain
unchanged; the cost model predicted the byte count exactly in all 42
distributed configurations measured. We additionally protect job submission
and result provenance with the NIST-standardised post-quantum primitives
ML-KEM-768 and ML-DSA-65, and measure that overhead: $1.49\%$ of the runtime
of the heaviest job in our suite, of which only $189\,\mu$s is lattice
arithmetic. All results are from one laptop; we are explicit throughout about
what that does and does not support.
\end{abstract}

\section{Introduction}

A pure quantum state of $n$ qubits is a unit vector
\begin{equation}
\ket{\psi} = \sum_{i=0}^{2^n-1} \alpha_i \ket{i}, \qquad
\sum_i |\alpha_i|^2 = 1 ,
\end{equation}
where we write $\ket{\cdot}$ for a computational basis state. Storing the
amplitudes explicitly costs $O(2^n)$: in double precision, 30 qubits is
16\,GiB and 40 qubits is 16\,TiB. Single-machine simulation therefore becomes
memory-bound long before it becomes compute-bound, and the standard response is
to distribute the state vector across the memory of many nodes
\cite{deraedt2007,quest2019}.

Distribution introduces a cost that does not exist on one machine.
With $P = 2^p$ ranks, each rank holds $2^{n-p}$ amplitudes and the high $p$
bit positions of a basis index select the rank that owns it. A gate acting on a
qubit in one of those positions may have to move amplitudes between ranks. The
assignment of logical qubits to bit positions is, however, free: nothing forces
qubit $q$ to occupy position $q$. The default assignment --- the identity
permutation --- is an accident of indexing, and for some circuits it is close
to the worst possible choice.

This report describes a system built around three questions.

\begin{description}
\item[RQ1] Can communication-aware qubit placement reduce MPI communication
  volume and wall-clock time for distributed state-vector simulation?
\item[RQ2] What overhead does authenticating and encrypting HPC quantum jobs
  and results with standardised post-quantum cryptography introduce?
\item[RQ3] How do small Grover and Shor experiments illustrate the
  cryptographic motivation for post-quantum cryptography?
\end{description}

The simulator is implemented from scratch in C++20 with MPI and OpenMP.
Qiskit \cite{qiskit} is used only as an external correctness oracle in the test
suite; it is never imported by the runtime.

\section{Background}

\subsection{State-vector simulation}

A single-qubit gate with matrix $U$ acting on qubit $q$ pairs each amplitude
$\alpha_i$ with $\alpha_{i \oplus 2^q}$ and applies $U$ to the pair. Enumerating
the $2^{n-1}$ pair representatives touches every amplitude exactly once, which
makes the kernel a strided, embarrassingly parallel sweep --- the reason
state-vector simulation threads well and the reason it is bandwidth-bound.

Two structural properties of a gate determine its behaviour under
partitioning. A \emph{diagonal} gate multiplies each amplitude by a scalar
that depends only on that amplitude's index, so it never moves data. A
\emph{control} qubit is only read, never written. Both properties turn out to
matter more than the gate's name.

\subsection{Distributed-memory simulation}

Established distributed simulators partition the state vector across
processes and exchange shards for gates acting on the high-order qubits
\cite{deraedt2007,quest2019,cuquantum}. The usual treatment of placement is to
apply gates to whichever qubits they name, occasionally reordering qubits so
that the expensive ones become local. We make that reordering the object of
study: the runtime carries an explicit logical-to-physical permutation, a cost
model predicts its consequences in bytes, and a search chooses it.

\subsection{Post-quantum cryptography}

NIST standardised module-lattice key encapsulation as ML-KEM (FIPS 203
\cite{fips203}) and module-lattice signatures as ML-DSA (FIPS 204
\cite{fips204}) in 2024. We use the reference implementations in liboqs
\cite{liboqs} through its Python binding, with HKDF-SHA256 \cite{rfc5869} for
key derivation and AES-256-GCM \cite{sp800038d} for payload encryption. No
cryptographic primitive is implemented in this project.

\section{System architecture}

AegisQ-HPC is four cooperating layers.

\begin{enumerate}
\item A \textbf{front end} providing a circuit IR over twelve gates
  (\texttt{x y z h s t rx ry rz cx cz swap}) and a documented OpenQASM subset.
  Anything outside the subset is rejected with a line number rather than
  ignored.
\item A \textbf{native core} in C++20: local kernels shared by single-process
  and distributed execution, an MPI runtime, and a communication profiler.
\item A \textbf{compiler} comprising an analytical communication cost model, a
  placement search and a gate-fusion pass.
\item A \textbf{security layer} providing ML-KEM/ML-DSA job envelopes and
  signed, Merkle-committed result provenance.
\end{enumerate}

A deliberately simple NumPy reference simulator is the correctness oracle for
everything else; the C++ and MPI engines are validated against it gate by gate
and on randomised circuits.

\section{Distributed simulation runtime}

\subsection{Partitioning}

With $P = 2^p$ ranks and $n$ qubits, a physical basis index splits as
\begin{equation}
\text{physical} = (\text{rank} \ll L) \mathbin{|} \text{local}, \qquad L = n - p ,
\end{equation}
so the low $L$ bit positions are \emph{local} and the high $p$ are
\emph{global}. A layout object carries a permutation
$\pi : [0,n) \to [0,n)$ mapping logical qubits to bit positions, defaulting to
the identity. Every kernel reads placement only through predicates
(\texttt{is\_local}, \texttt{global\_position}, \texttt{partner\_rank}), so
changing the permutation requires no change to any kernel.

The initial state $\ket{0\cdots 0}$ is physical index $0$ under every
permutation and lives on rank $0$, so initialisation needs no communication.

\subsection{Which gates communicate}

Let $S = 2^{L}$ be the shard size in amplitudes and $w$ the bytes per
amplitude. Table~\ref{tab:costs} summarises the traffic per gate placement;
it is the cost model and it is also exactly what the runtime does.

\begin{table}[t]
\centering
\begin{tabular}{llrr}
\toprule
Gate & Placement & Bytes (all ranks) & Messages \\
\midrule
\texttt{z s t rz cz} & any & $0$ & $0$ \\
\texttt{x y h rx ry} & local & $0$ & $0$ \\
\texttt{x y h rx ry} & global & $P S w$ & $P$ \\
\texttt{cx} & local target & $0$ & $0$ \\
\texttt{cx} & local control, global target & $P S w / 2$ & $P$ \\
\texttt{cx} & global control, global target & $P S w / 2$ & $P/2$ \\
\texttt{swap} & both local & $0$ & $0$ \\
\texttt{swap} & one global & $P S w / 2$ & $P$ \\
\texttt{swap} & both global & $P S w / 2$ & $P/2$ \\
\bottomrule
\end{tabular}
\caption{Communication per gate and placement.}
\label{tab:costs}
\end{table}

Three entries deserve comment.

\emph{Diagonal gates are free even on a global qubit.} The rank identifier
fixes that qubit's value for every amplitude the rank holds, so the gate
collapses to a single scalar multiply of the whole shard.

\emph{A global control is cheap; a global target is not.} A control is read,
so a rank that does not satisfy it skips the gate entirely and a rank that does
proceeds locally. A target is written, and writing a global bit moves
amplitudes.

\emph{Half-shard exchanges are real, not an approximation.} In a \texttt{cx}
with a local control, only amplitudes whose control bit is set participate.
Those form contiguous runs of length $2^{\pi(c)}$, so they are packed with a
strided block copy, exchanged, and scattered back --- half the bytes of a whole
shard exchange.

All pairwise traffic uses a symmetric \texttt{MPI\_Sendrecv}, which cannot
deadlock, into a separate buffer, so no amplitude is overwritten before it is
consumed. Buffers are reused across gates and transfers are chunked at $2^{28}$
elements because MPI element counts are \texttt{int}-typed.

\subsection{Measurement}

Shot sampling draws \texttt{shots} uniform values in $[0, 1)$ from a seeded
Mersenne Twister, sorts them once, and walks the local amplitudes accumulating
probability mass; each rank claims only the draws that fall inside its own
interval, whose offset comes from an \texttt{MPI\_Exscan}. Because the draw
sequence depends only on the seed and the shot count, a circuit sampled on 8
ranks produces \emph{identical} counts to the same circuit sampled in one
process --- a property asserted in the test suite and relied on by the
benchmark methodology.

\subsection{Instrumentation}

Every transfer the runtime issues is timed and counted inside the function that
issues it: sends, receives, pairwise exchanges, bytes in each direction,
collectives, and a per-opcode breakdown. What is counted is the payload handed
to MPI; what is not counted is MPI's own protocol overhead and the internal
traffic of collectives. Reported volumes are therefore a lower bound on wire
traffic and an exact account of what the algorithm requested --- which is the
quantity placement controls.

\section{Communication-aware qubit placement}

Because byte volume depends only on \emph{which} qubits are global and not on
their ordering within each group, the search space is the $\binom{n}{p}$
subsets. For 24 qubits on 8 ranks that is 2024 candidates, so an exhaustive
search is affordable and is optimal with respect to the model; beyond a
configurable budget the search falls back to a greedy start with pairwise local
improvement and reports that it is not guaranteed optimal.

The objective is lexicographic: bytes first, message count as a tiebreak.
Message count distinguishes placements that move identical volume in different
numbers of transfers, which bytes alone cannot.

The model is not trusted on its own. For every circuit in the evaluation, the
predicted byte count is compared against the profiler's measurement on the live
runtime, including under permuted placements and in single precision.

\section{Post-quantum job protocol}

A job travels as one file so it can move over whatever transport a cluster
already has. The client encapsulates to the cluster's ML-KEM public key,
derives an AES key with HKDF bound to the protocol version and the job
identifier, encrypts the manifest and the circuit \emph{together} with
AES-256-GCM using the public header as associated data, and signs the header
and ciphertext with ML-DSA.

The double binding is deliberate: editing a public field breaks decryption as
well as the signature, so a job can never run under parameters other than those
that were signed. Everything signed or hashed is serialised canonically
(sorted keys, no insignificant whitespace), and the parser re-serialises what it
parsed and compares bytes, rejecting input that was not already canonical.

On receipt the cluster checks, in order: structure; the signing key against a
local allow-list; the ML-DSA signature; the addressee; replay state; then
decapsulation, decryption and the circuit hash. Nothing inside the envelope is
trusted before the signature verifies.

Results are committed with a Merkle tree following RFC 6962 \cite{rfc6962}:
leaves and internal nodes are domain-separated, and odd leaf counts split at the
largest power of two rather than duplicating the last leaf, avoiding the
ambiguity in which two distinct trees share a root. The manifest --- artefact
digests, environment, placement, timings --- is signed with the cluster's
ML-DSA key.

We state the resulting claim precisely. A valid signature means the record was
produced by a holder of the cluster's secret key and has not been modified. It
does \emph{not} mean the computation was performed correctly: a compromised
node signs a wrong answer with a valid key. Verifying provenance is not
verifying computation.

\section{Evaluation}

\subsection{Setup}

All measurements were taken on one Apple M2 laptop (8 logical cores, 16\,GiB),
macOS 15.6, Apple Clang 17, Open MPI 5.0.10, with shared-memory MPI and no
inter-node network. Each configuration runs a discarded warm-up repeat followed
by three timed repeats; we report the minimum wall time with the median and
spread recorded alongside. Thread policy is an explicit variable: under
\emph{one thread per rank} the total core count grows with the rank count
(classical scaling), while under \emph{fixed total cores} $\text{ranks} \times
\text{threads}$ is held at the core count, which measures the cost of
partitioning rather than the benefit of adding processors. The two are never
combined in one table.

\subsection{Correctness}

The native and distributed engines are validated against the NumPy reference
simulator per gate, per placement case, and on randomised circuits at 1, 2, 4
and 8 ranks; the reference is in turn cross-validated against Qiskit on 100
seeded random circuits per run. The full suite is 666 single-process tests plus
109 distributed tests per rank count.

\subsection{RQ1: communication-aware placement}

% Generated by scripts/generate_report.py -- do not edit.
\begin{table}[t]
\centering
\small
\begin{tabular}{lrrrrrr}
\toprule
 & \multicolumn{3}{c}{MPI bytes sent} & \multicolumn{3}{c}{Wall time} \\
\cmidrule(lr){2-4}\cmidrule(lr){5-7}
Circuit & Default & Optimized & Reduction & Default & Optimized & Change \\
 & (MiB) & (MiB) & & (ms) & (ms) & \\
\midrule
qft & 936 & 120 & 87.2\% & 358 & 283 & -20.9\% \\
grover & 1792 & 496 & 72.3\% & 566 & 444 & -21.6\% \\
random & 392 & 240 & 38.8\% & 166 & 129 & -22.1\% \\
ising & 432 & 368 & 14.8\% & 164 & 160 & -2.6\% \\
ghz & 24 & 24 & 0.0\% & 10 & 16 & +58.7\% \\
\bottomrule
\end{tabular}
\caption{Measured MPI traffic and wall time under the default and communication-aware placements, 20 qubits on 8 ranks.}
\label{tab:mapping}
\end{table}


Table~\ref{tab:mapping} gives the measured effect at 8 ranks.
The quantum Fourier transform benefits most: in the textbook decomposition
every controlled phase targets the higher-indexed qubit, and the identity
placement puts exactly those qubits on the global positions. Moving them local
removes $87.2\%$ of the traffic and $24.3\%$ of the wall time. Grover's dense
multi-controlled blocks behave similarly. The GHZ chain shows no reduction at
all --- every qubit but one is a \texttt{cx} target exactly once, so no
placement is better than another. We report that as a result: a heuristic that
claimed a win on every circuit would be the suspicious one.

Wall time does not track bytes proportionally. The Ising chain sheds $14.8\%$
of its traffic with no measurable time improvement, which is what one expects
when communication is not the bottleneck --- here, shared-memory MPI on one
socket, where a ``transfer'' is a memory copy.

The cost model predicted the byte count and the message count \emph{exactly} in
all 42 distributed configurations measured.

\subsection{A second lever: gate fusion}

Placement decides which gates communicate; it cannot reduce how often a given
qubit is acted on. A run of consecutive single-qubit gates on a global qubit
costs one whole-shard exchange \emph{per gate}, although the run's product is
a single $2\times2$ matrix. Multiplying the run out before execution replaces
those exchanges with one.

We carry the fused matrix in full rather than as Euler angles. There is no
hardware target to compile for, and the explicit form is exact including the
global phase, so fused and unfused circuits are compared amplitude by
amplitude rather than through a fidelity. Two further properties are checked
by property-based tests over arbitrary circuits: fusion never increases
predicted communication, and a run of diagonal gates fuses into a diagonal
matrix that the runtime and the cost model both still recognise as free --- a
fused gate is classified by inspecting its matrix, not its opcode.

% Generated by scripts/generate_report.py -- do not edit.
\begin{table}[t]
\centering
\small
\begin{tabular}{lrrrr}
\toprule
Circuit & Baseline (MiB) & Fusion only & Placement only & Both \\
\midrule
qft & 936 & 0.0\% & 87.2\% & 87.2\% \\
grover & 1792 & 16.1\% & 72.3\% & 86.6\% \\
random & 392 & 8.2\% & 38.8\% & 53.1\% \\
ising & 432 & 0.0\% & 14.8\% & 14.8\% \\
ghz & 24 & 0.0\% & 0.0\% & 0.0\% \\
\bottomrule
\end{tabular}
\caption{Measured reduction in MPI traffic from each optimisation lever at 8 ranks, 20 qubits.}
\label{tab:levers}
\end{table}


Table~\ref{tab:levers} gives the measured $2\times2$. The levers are not
additive. For random circuits fusion alone removes $8.2\%$ of traffic and
placement alone $38.8\%$, while the combination removes $53.1\%$: fusing
changes the cost landscape, so the placement search that follows finds a
different and better assignment. For the QFT, whose single-qubit runs are
already short, fusion contributes nothing and placement does all the work.

% Generated by scripts/generate_report.py -- do not edit.
\begin{table}[t]
\centering
\small
\begin{tabular}{lrrrrr}
\toprule
Circuit & Qubits & Ranks & Wall (s) & Speedup & Efficiency \\
\midrule
ising & 22 & 1 & 1.815 & 1.00 & 100\% \\
ising & 22 & 2 & 0.961 & 1.89 & 94\% \\
ising & 22 & 4 & 0.681 & 2.66 & 67\% \\
ising & 22 & 8 & 0.750 & 2.42 & 30\% \\
qft & 22 & 1 & 5.211 & 1.00 & 100\% \\
qft & 22 & 2 & 2.765 & 1.88 & 94\% \\
qft & 22 & 4 & 1.907 & 2.73 & 68\% \\
qft & 22 & 8 & 1.988 & 2.62 & 33\% \\
\bottomrule
\end{tabular}
\caption{Strong scaling with one thread per rank.}
\label{tab:strong}
\end{table}


Table~\ref{tab:strong} gives strong scaling with one thread per rank. Efficiency
falls from $92\%$ at 2 ranks to about $30\%$ at 8, on a machine with 4
performance and 4 efficiency cores and no separate interconnect. Under the
fixed-total-cores policy the same QFT reaches only $0.87\times$ --- that is, on
constant hardware, partitioning costs more than it saves, which is the honest
framing of what MPI buys on a single node.

\subsection{RQ2: cost of the post-quantum layer}

% Generated by scripts/generate_report.py -- do not edit.
\begin{table}[t]
\centering
\small
\begin{tabular}{llrr}
\toprule
Algorithm & Operation & Median ($\mu$s) & Bytes \\
\midrule
ML-DSA-44 & keygen & 33.0 & 1312 \\
ML-DSA-44 & sign & 63.0 & 2420 \\
ML-DSA-44 & verify & 31.3 & 2420 \\
ML-DSA-65 & keygen & 53.8 & 1952 \\
ML-DSA-65 & sign & 91.0 & 3309 \\
ML-DSA-65 & verify & 50.7 & 3309 \\
ML-DSA-87 & keygen & 89.4 & 2592 \\
ML-DSA-87 & sign & 139.5 & 4627 \\
ML-DSA-87 & verify & 83.4 & 4627 \\
ML-KEM-1024 & decapsulate & 29.4 & 32 \\
ML-KEM-1024 & encapsulate & 25.7 & 1568 \\
ML-KEM-1024 & keygen & 24.5 & 1568 \\
ML-KEM-512 & decapsulate & 14.5 & 32 \\
ML-KEM-512 & encapsulate & 12.8 & 768 \\
ML-KEM-512 & keygen & 12.1 & 800 \\
ML-KEM-768 & decapsulate & 21.2 & 32 \\
ML-KEM-768 & encapsulate & 18.6 & 1088 \\
ML-KEM-768 & keygen & 16.8 & 1184 \\
\bottomrule
\end{tabular}
\caption{Post-quantum primitive cost and artefact sizes.}
\label{tab:pqc}
\end{table}


Table~\ref{tab:pqc} gives primitive costs. End to end, packing a 20-qubit job
bundle takes $3.1$\,ms and verifying, decrypting and opening it $5.6$\,ms, with
$7083$ bytes of envelope overhead independent of circuit size. Against the
heaviest job in our suite ($582$\,ms, $1792$\,MiB of MPI traffic) that is
$1.49\%$ of runtime.

The decomposition matters more than the total: only $189\,\mu$s of that is
lattice arithmetic. The remainder is canonical JSON serialisation and base64
encoding of the circuit. A reader who concluded ``post-quantum cryptography is
expensive'' from the end-to-end number would be drawing the wrong lesson; the
serialisation is what an optimisation would need to target.

\subsection{RQ3: why post-quantum cryptography}

% Generated by scripts/generate_report.py -- do not edit.
\begin{table}[t]
\centering
\small
\begin{tabular}{rrrrr}
\toprule
$N$ & Qubits & Classical queries & Grover queries & Measured success \\
\midrule
4 & 2 & 2.5 & 1 & 100.0\% \\
8 & 4 & 4.5 & 2 & 95.4\% \\
16 & 6 & 8.5 & 3 & 96.6\% \\
32 & 8 & 16.5 & 4 & 99.9\% \\
64 & 10 & 32.5 & 6 & 99.7\% \\
128 & 12 & 64.5 & 8 & 99.6\% \\
256 & 14 & 128.5 & 12 & 100.0\% \\
\bottomrule
\end{tabular}
\caption{Grover oracle queries against classical search, measured by simulation.}
\label{tab:grover}
\end{table}


Table~\ref{tab:grover} reports simulated Grover searches. Every row is a run,
not a formula: the measured probability of finding the marked state agrees with
the amplitude-amplification prediction to within a few tenths of a percent, and
$256$ items are searched in 12 oracle queries against a classical expectation
of $128.5$.

That quadratic factor \cite{grover1996} is precisely why post-quantum guidance
doubles symmetric key sizes rather than abandoning symmetric cryptography.
Shor's exponential advantage \cite{shor1997} is a different matter, and our
implementation illustrates it at a size that makes the point without
overstating it: the order-finding circuit factors $15$ and $21$. Controlled
modular multiplication is synthesised generically --- the map is a permutation
of basis states, decomposed into transpositions, each realised by a CNOT ladder
and one multi-controlled $X$ --- which is exponential in the register width and
therefore stops where the simulation does.

\section{Limitations}

Every measurement here comes from a single laptop with shared-memory MPI. There
is no inter-node network, so communication times are best-case; on a real
interconnect the placement optimisation should matter \emph{more}, but that is
a prediction, not a result. Problem sizes are 20--23 qubits. There is no GPU
backend and no GPU was available. We make no performance comparison against
other simulators, because a fair comparison requires care about which
optimisations each applies.

The simulator is exact and noiseless, supports twelve gates and terminal
measurement only, and requires power-of-two rank counts. The cost model counts
bytes and messages and does not model topology, congestion, or overlap of
computation with communication. Placement is chosen once per circuit; dynamic
re-mapping between circuit segments is future work.

The security layer is a research prototype. A malicious cluster administrator,
memory disclosure during execution, side channels (including the communication
volume this project deliberately exposes), denial of service, key distribution
and revocation are all outside the threat model.

\section{Future work}

Three directions follow directly from what is measured here. \emph{Dynamic
re-mapping}: splitting a circuit into windows and re-permuting when the
predicted saving exceeds the permutation cost, which the cost model already has
the machinery to decide. \emph{Cluster measurements}: the same experiments on
multiple nodes, where communication is a network rather than a memory copy and
the byte reductions reported here should translate into larger time reductions.
\emph{Topology awareness}: extending the cost model beyond byte counts to
account for the interconnect a job is actually running on. The measured
interaction between placement and fusion also suggests that the two should be
searched jointly rather than applied in sequence.

\section{Conclusion}

Partitioning a quantum state vector forces a choice that is usually made by
accident. Making it deliberately --- with a cost model in bytes, a search over
placements, and instrumentation that measures what actually crossed the network
--- removed up to $87\%$ of MPI traffic on the circuits we measured, and the
model predicted the byte counts exactly in every distributed configuration.
Protecting the same jobs with NIST-standardised post-quantum cryptography cost
about $1.5\%$ of runtime, dominated by serialisation rather than lattice
arithmetic. Both numbers come from one laptop, and the repository says so in
every place they appear.

\bibliographystyle{plain}
\bibliography{references}

\end{document}
